cent vingt ans... ». On ne connaît pas dans les annales d'Alban un autre
cas de longévité semblable. Un descendant de cette famille :
Antoine Pujol fera, plus d'un siècle après, graver son nom dans la
pierre. On peut encore l'y lire aujourd'hui, au-dessus d'une lucarne
qui donne le jour à ce qui est maintenant une étable. Cet Antoine
Pujol pourrait bien être celui que le cadastre de 1693 mentionne
comme « lieutenant en la justice du vicomté de Paulin et de la
1\iere compagnie de bourgeoisie du lieu de le Tapie\footnote{La Tapie?} ».

Le nouveau seigneur d'Alban : Jean Albert, eut bientôt maille à
partir avec les protestants. Ceux-ci occupent Paulin. Ils sont aussi
à Trébas. En 1567, ils se sont emparé d'Alban où on les trouve
encore en 1568. En effet, l'évêque d'Albi mentionne à cette époque
parmi les lieux occupés par les réformés : « Albaing (Alban),
Trébas, Mules (Miolles), St Jean de Janes... Teillet... »
L'occupation d'Alban paraît d'assez courte durée puisque en Juin
1570, une garnison de quinze hommes sous le commandement du
capitaine de Nougaret a été placée à Alban par les soins de
Villeneuve, lieutenant du duc de Joyeuse. En quittant notre cité,
nos calvinistes ont emmené les deux consuls d'Alban. La cité est
à ce point ruinée qu'elle ne peut payer en 1571 la rançon de leur
retour (octobre 1571).